% =========================================================
% TOPIC 2 -- PROOFS
% =========================================================

\section{Proofs}

% =========================================================
% TOPIC TITLE SLIDE
% =========================================================

\begin{frame}{Proofs}

\begin{center}

\vspace{1cm}

{\Huge\textbf{Proofs}}

\vspace{0.6cm}

\textit{Constructing, evaluating, and communicating mathematical proofs}

\end{center}

\end{frame}

% =========================================================
% QUESTION 1 -- PART 1
% =========================================================

\begin{frame}{Proofs -- Question 1}

\begin{block}{Algorithm Correctness}

A programmer claims that the following algorithm correctly determines
whether an integer $n$ is even:

\[
\text{If } n \bmod 2 = 0,\text{ return ``Even'';}
\]

\[
\text{otherwise, return ``Odd''.}
\]

The programmer wants to justify that the algorithm is correct for every
integer input.

\end{block}

\end{frame}

% =========================================================
% QUESTION 1 -- PART 2
% =========================================================

\begin{frame}{Proofs -- Question 1 (Continued)}

\begin{block}{Question}

Construct a mathematical proof that the algorithm correctly classifies
every integer as either even or odd.

\medskip

Your proof should explain why every integer $n$ must fall into one of the
two cases represented by the algorithm.

\end{block}

\end{frame}

% =========================================================
% ANSWER 1 -- PART 1
% =========================================================

\begin{frame}{Proofs -- Question 1: Answer}

\begin{exampleblock}{Proof}

Let $n$ be any integer.

By the division algorithm, when $n$ is divided by $2$, the remainder must
be either $0$ or $1$.

Therefore,

\[
n=2q
\]

for some integer $q$, or

\[
n=2q+1
\]

for some integer $q$.

\end{exampleblock}

\end{frame}

% =========================================================
% ANSWER 1 -- PART 2A
% =========================================================

\begin{frame}{Proofs -- Question 1: Answer (Continued)}

\begin{exampleblock}{Proof}

If
\[
n=2q,
\]

then
\[
n\bmod 2=0,
\]

so the algorithm returns ``Even''.

\medskip

If
\[
n=2q+1,
\]

then
\[
n\bmod 2=1,
\]

so the algorithm returns ``Odd''.

\end{exampleblock}

\end{frame}


% =========================================================
% QUESTION 2 -- PART 1
% =========================================================

\begin{frame}{Proofs -- Question 2}

\begin{block}{Data Validation Rule}

A data-processing system stores integer values that are required to be
multiples of $3$.

A developer claims:

\begin{center}
\textit{``The sum of two valid stored values is also a valid stored value.''}
\end{center}

Suppose $a$ and $b$ are two valid stored values.

\end{block}

\end{frame}

% =========================================================
% QUESTION 2 -- PART 2
% =========================================================

\begin{frame}{Proofs -- Question 2 (Continued)}

\begin{block}{Question}

Prove the developer's claim mathematically.

\medskip

Your proof should begin by expressing $a$ and $b$ as multiples of $3$ and
show that their sum is also a multiple of $3$.

\end{block}

\end{frame}

% =========================================================
% ANSWER 2 -- PART 1
% =========================================================

\begin{frame}{Proofs -- Question 2: Answer}

\begin{exampleblock}{Proof}

Since $a$ and $b$ are valid stored values, both are multiples of $3$.

Therefore, there exist integers $m$ and $n$ such that

\[
a=3m
\]

and

\[
b=3n.
\]

Consider their sum:

\[
a+b=3m+3n.
\]

\end{exampleblock}

\end{frame}

% =========================================================
% ANSWER 2 -- PART 2
% =========================================================

\begin{frame}{Proofs -- Question 2: Answer (Continued)}

\begin{exampleblock}{Proof}

Factor out $3$:

\[
a+b=3(m+n).
\]

Since $m+n$ is an integer,

\[
a+b
\]

is a multiple of $3$.

Therefore, the sum of two valid stored values is also valid.

\[
\boxed{a+b\text{ is a multiple of }3}
\]

Hence, the developer's claim is proved.

\end{exampleblock}

\end{frame}

% =========================================================
% QUESTION 3 -- PART 1
% =========================================================

\begin{frame}{Proofs -- Question 3}

\begin{block}{Network Identifier Property}

A network system assigns integer identifiers to devices.

A particular subsystem uses only identifiers of the form

\[
2k+1,
\]

where $k$ is an integer.

An engineer observes that combining two such identifiers by addition
always produces an even number.

\end{block}

\end{frame}

% =========================================================
% QUESTION 3 -- PART 2
% =========================================================

\begin{frame}{Proofs -- Question 3 (Continued)}

\begin{block}{Question}

Prove that the sum of any two identifiers of the form $2k+1$ is even.

\medskip

Let the two identifiers be

\[
a=2m+1
\]

and

\[
b=2n+1,
\]

where $m,n\in\mathbb{Z}$.

Show that $a+b$ can be written in the form $2r$ for some integer $r$.

\end{block}

\end{frame}

% =========================================================
% ANSWER 3 -- PART 1
% =========================================================

\begin{frame}{Proofs -- Question 3: Answer}
\vspace{-0.5cm}
\begin{exampleblock}{Proof}

Let
\[
a=2m+1
\]

and
\[
b=2n+1,
\]

where
\[
m,n\in\mathbb{Z}.
\]

Their sum is
\[
a+b=(2m+1)+(2n+1).
\]

Simplifying,
\[
a+b=2m+2n+2.
\]

\end{exampleblock}

\end{frame}

% =========================================================
% ANSWER 3 -- PART 2
% =========================================================

\begin{frame}{Proofs -- Question 3: Answer (Continued)}

\begin{exampleblock}{Proof}

Factor out $2$:
\[
a+b=2(m+n+1).
\]

Since $m$, $n$, and $1$ are integers,
\[
m+n+1\in\mathbb{Z}.
\]

Let
\[
r=m+n+1.
\]

Then
\[
a+b=2r.
\]

Therefore, $a+b$ is even.
\[
\boxed{\text{The sum of two odd integers is even.}}
\]

\end{exampleblock}

\end{frame}

% =========================================================
% QUESTION 4 -- PART 1
% =========================================================

\begin{frame}{Proofs -- Question 4}

\begin{block}{Software Access Rule}

A software system assigns each user an integer access score.

The security team claims:

\begin{center}
\textit{``If a user's access score is divisible by $4$, then it is even.''}
\end{center}

A developer attempts to prove this by saying:

\begin{center}
\textit{``All even numbers are divisible by $4$, so any number divisible
by $4$ is even.''}
\end{center}

\end{block}

\end{frame}

% =========================================================
% QUESTION 4 -- PART 2
% =========================================================

\begin{frame}{Proofs -- Question 4 (Continued)}

\begin{block}{Question}

Determine whether the developer's reasoning is a valid proof.

\medskip

If it is not valid, identify the error in the reasoning and provide a
correct proof of the security team's claim.

\end{block}

\end{frame}

% =========================================================
% ANSWER 4 -- PART 1
% =========================================================

\begin{frame}{Proofs -- Question 4: Answer}

\begin{exampleblock}{Analysis}

The developer's reasoning is not valid.

The statement

\[
\text{``All even numbers are divisible by }4\text{''}
\]

is false.

For example,

\[
2
\]

is even, but $2$ is not divisible by $4$.

Therefore, the starting statement used by the developer is incorrect.

\end{exampleblock}

\end{frame}

% =========================================================
% ANSWER 4 -- PART 2
% =========================================================

\begin{frame}{Proofs -- Question 4: Answer (Continued)}
\vspace{-0.5cm}
\begin{exampleblock}{Correct Proof}

Let $n$ be an integer divisible by $4$.

Then there exists an integer $k$ such that
\[
n=4k.
\]
Since
\[
4k=2(2k),
\]

we can write
\[
n=2(2k).
\]

Because $2k$ is an integer, $n$ is even.

Therefore,
\[
\boxed{4\mid n\Rightarrow 2\mid n}.
\]

Hence, the security team's claim is true.

\end{exampleblock}

\end{frame}

% =========================================================
% QUESTION 5 -- PART 1
% =========================================================

\begin{frame}{Proofs -- Question 5}

\begin{block}{Hashing System Property}

A hashing system maps an integer input $n$ to

\[
h(n)=n^2+n.
\]

A developer claims that the hash value $h(n)$ is always even, regardless
of the integer input.

\medskip

For example,

\[
h(3)=3^2+3=12.
\]

\end{block}

\end{frame}

% =========================================================
% QUESTION 5 -- PART 2
% =========================================================

\begin{frame}{Proofs -- Question 5 (Continued)}

\begin{block}{Question}

Prove that

\[
n^2+n
\]

is even for every integer $n$.

\medskip

Your proof should consider the possible parity of $n$ and show that the
result is even in every case.

\end{block}

\end{frame}

% =========================================================
% ANSWER 5 -- PART 1
% =========================================================

\begin{frame}{Proofs -- Question 5: Answer}

\begin{exampleblock}{Proof}

Let $n$ be any integer.

Every integer is either even or odd.

\medskip

\textbf{Case 1: $n$ is even}

Then there exists an integer $k$ such that

\[
n=2k.
\]

Therefore,

\[
n^2+n
=
(2k)^2+2k
=
4k^2+2k.
\]

\end{exampleblock}

\end{frame}

% =========================================================
% ANSWER 5 -- PART 2
% =========================================================

\begin{frame}{Proofs -- Question 5: Answer (Continued)}
\vspace{-0.5cm}
\begin{exampleblock}{Proof}

Factor out $2$:
\[
n^2+n
=
2(2k^2+k).
\]

Therefore, $n^2+n$ is even.

\medskip

\textbf{Case 2: $n$ is odd}

Let
\[
n=2k+1.
\]

Then
\[
n^2+n
=
(2k+1)^2+(2k+1)
\]
\[
=4k^2+6k+2
\]
\[
=2(2k^2+3k+1).
\]

Therefore, $n^2+n$ is even.

\end{exampleblock}

\end{frame}

% =========================================================
% ANSWER 5 -- PART 3
% =========================================================

\begin{frame}{Proofs -- Question 5: Answer (Final)}

\begin{exampleblock}{Conclusion}

In both possible cases,

\[
n^2+n
\]

is even.

Therefore,

\[
\boxed{n^2+n\text{ is even for every }n\in\mathbb{Z}.}
\]

Hence, the developer's claim about the hashing function is correct.

\medskip

\textbf{Proof idea:}

The argument uses a proof by cases based on whether the integer is even
or odd.

\end{exampleblock}

\end{frame}

% =========================================================
% END OF TOPIC 2
% =========================================================
